\documentclass[../main.tex]{subfiles}

\begin{document}
% -----------------------------
\section{Constant folding}
% -----------------------------
Now that everything is arith, we can start doing real optimization passes:
\begin{itemize}
    \item Fold arith.addi/muli/subi when both operands are constants
    \item Simplify relu(constant) to a constant
    \item Remove dead constants no longer used
\end{itemize}

What we want to do as a result is a pass that turns \texttt{relu((1+2)*3-10)}
into a single arith.constant 0 : i32 (or the computed value if positive).

Change import section:
\begin{lstlisting} [language=Python, caption={Importing SSAValue}]
from xdsl.ir import Region, Block, Operation, SSAValue
\end{lstlisting}
Delete \path{apply_lowering} function and replace it with:
\begin{lstlisting}[language=Python, caption={Apply pass implementation}]
def apply_pass(module: ModuleOp, pattern: RewritePattern) -> None:
    applier = GreedyRewritePatternApplier([pattern()])

    # Some xdsl builds have extra kwargs; enable recursion if available.
    init_sig = inspect.signature(PatternRewriteWalker.__init__)
    kwargs = {}
    for k in ("walk_regions", "apply_recursively", "walk_into_regions"):
        if k in init_sig.parameters:
            kwargs[k] = True

    walker = PatternRewriteWalker(applier, **kwargs)

    # Walk/rewrite the function bodies explicitly
    for top_block in module.body.blocks:
        for op in top_block.ops:
            if isinstance(op, func.FuncOp):
                if hasattr(walker, "rewrite_region"):
                    walker.rewrite_region(op.body)
                elif hasattr(walker, "rewrite_op"):
                    walker.rewrite_op(op)
                else:
                    raise RuntimeError("No suitable rewrite entrypoint found on PatternRewriteWalker.")
\end{lstlisting}
since we are going to reuse it for constant folding.
\\Add a couple of helper functions to deal with IntAttr and extract underlying integer:
\begin{lstlisting} [language=Python, caption={Helper function}]
# ---------------------------------------
# Helper functions
# ---------------------------------------

def _defining_op(val: SSAValue) -> Operation | None:
    # SSAValue in xdsl typically has .owner (def op)
    return getattr(val, "owner", None)

def _get_const_int_binop_args(op: Operation):
    """Return (a, b) if both operands are integer constants, else None."""
    if len(op.operands) != 2:
        return None

    lhs, rhs = op.operands
    lhs_op = _defining_op(lhs)
    rhs_op = _defining_op(rhs)
    if lhs_op is None or rhs_op is None:
        return None

    a = _int_from_constant_op(lhs_op)
    b = _int_from_constant_op(rhs_op)
    if a is None or b is None:
        return None

    return a, b

def _int_from_constant_op(op: Operation) -> int | None:
    """
    Extract integer from arith.constant.
    Returns None if not an integer constant we can read.
    """
    if op.name != "arith.constant":
        return None

    # In many versions: arith.ConstantOp has attribute `.value`
    v = getattr(op, "value", None)
    if v is None:
        return None

    # Common shapes: IntegerAttr(value=...), or has `.value`/`.data`
    for attr_name in ("value", "data"):
        if hasattr(v, attr_name):
            try:
                return int(getattr(v, attr_name))
            except Exception:
                pass

    # Last resort: try to parse int from string like "1 : i32"
    # (kept minimal; only if the above fails)
    try:
        s = str(op)
        # crude parse: find "arith.constant " then grab next token
        if "arith.constant" in s:
            tok = s.split("arith.constant", 1)[1].strip().split()[0]
            return int(tok)
    except Exception:
        pass

    return None


def _result_width(op: Operation, default: int = 32) -> int:
    # Try to read integer width from result type if available
    if not op.results:
        return default

    ty = op.results[0].type
    w = getattr(ty, "width", None)
    if w is None:
        return default

    # In our xdsl, width is an IntAttr, not a Python int
    for attr_name in ("data", "value"):
        if hasattr(w, attr_name):
            return int(getattr(w, attr_name))

    # fallback: try direct int conversion if supported
    try:
        return int(w)
    except Exception:
        return default
\end{lstlisting}

\begin{lstlisting} [language=Python, caption={Constant folding pattern implementation}]
# --------------------------------------------
# Constant folding pattern
# --------------------------------------------
class FoldArithInts(RewritePattern):
    def match_and_rewrite(self, op: Operation, rewriter: PatternRewriter):
        # Special-case: fold relu lowering if it became maxsi
        if op.name == "arith.maxsi":
            args = _get_const_int_binop_args(op)
            if args is None:
                return
            a, b = args
            width = _result_width(op, 32)
            c = arith.ConstantOp.from_int_and_width(max(a, b), width)
            rewriter.replace_op(op, new_ops=[c], new_results=[c.result], safe_erase=True)
            return

        # Fold binary integer ops when both operands are integer constants
        if op.name not in ("arith.addi", "arith.muli", "arith.subi"):
            return

        args = _get_const_int_binop_args(op)
        if args is None:
            return
        a, b = args

        res_by_op = {
            "arith.addi": a + b,
            "arith.muli": a * b,
            "arith.subi": a - b,
        }
        res = res_by_op[op.name]

        width = _result_width(op, 32)
        c = arith.ConstantOp.from_int_and_width(res, width)
        rewriter.replace_op(op, new_ops=[c], new_results=[c.result], safe_erase=True)
\end{lstlisting}
Rename LowerHCAddPattern class to LowerHCPattern, since it now contains more operations.
\\Add constant folding pass in the main function:
\begin{lstlisting} [language=Python, caption={Main function}]
def main():
    ctx = build_context()
    module = build_module(ctx)

    print("=== HIGH LEVEL (hc.*) ===")
    print(module)

    # Apply lowering
    apply_pass(module, LowerHCPattern)

    print("\n=== AFTER LOWERING  (arith.*) ===")
    print(module)

    # Apply constant folding
    apply_pass(module, FoldArithInts)
    print("\n=== AFTER CONSTANT FOLDING ===")
    print(module)
\end{lstlisting}

\end{document}
